% Notes of the Montreal seminar, 1970 (cote 129, typed part): archive cote 129, site page 30 (PDF page 31).
% typescript. Draft compared: none.
% AI-transcribed by Claude Opus 5.5 (claude-opus-5-5, adaptive thinking, effort high), Message Batches API, batch mtl129_full (msgbatch_01AGnDKHwy6ATCM7ehUa6263), 2026-11-13; not yet checked by eye.
\scanpage{129/30}
\mspage{\hand{II.2}}

de $A \otimes_{\underline{O}_S} k(s)$ . L'application $s \to$ rang de $G$ en $s$ est une application
localement constante sur $S$ . Si ce rang est constant égal à
$n$ , on dit que $G$ est un \uline{groupe de rang} $n$ \uline{sur} $S$ .

1.3. \uline{Exemples de groupes finis localement libres}

1.3.1. Soit $G$ un groupe (abstrait) d'ordre $n$ . Soit $\underline{G}_S$ le fais-%
ceau constant sur $S$ défini par
\[ \underline{G}_S(S') = G \]
pour tout S-schéma $S'$ . $\underline{G}_S$ est un groupe libre de rang $n$ sur $S$ ; en
effet son algèbre affine est isomorphe à $\underline{O}_S^n$ .

1.3.2. Soit $\mathbb{G}_{m|S}$ le groupe multiplicatif sur $S$ , défini par
\[ \mathbb{G}_{m|S}(S') = \Gamma(S', \underline{O}_{S'})^* \]
pour tout S-schéma $S'$ . Pour tout $n \geq 1$ , on a l'isomorphisme
\[ n\ \mathrm{Id} : \mathbb{G}_{m|S} \to \mathbb{G}_{m|S} \qquad x \to x^n . \]
\margin{$\times$ \ $\mathbb{\mu}_{n|S}$}Par définition, $\deleted{\mu}_{n|S}$ , \uline{groupe des racines} $n^{\text{ièmes}}$ \uline{de l'unité sur} $S$ , est
le noyau de cet homomorphisme. C'est un groupe libre fini derang $n$ sur $S$ ;
en effet, son algèbre affine est isomorphe à $\underline{O}_S[T]/(T^n-1)$ .

1.3.3. Soit $\mathbb{G}_{a|S}$ le groupe additif sur $S$ , défini par
\[ \mathbb{G}_{a|S}(S') = \Gamma(S', \underline{O}_{S'}) \]
\margin{$\times$ \ $\mathbb{F}_p$}pour tout S-schéma $S'$ . Si $S$ est un schéma sur le corps premier $\mathbb{F}_p$ , on
a un homomorphisme de Frobenius

\margin{30}
